% Completed supersingular calculation; independently audited by root and fhi.
\subsection*{D4s: the supersingular Igusa hull and the canonical elliptic subgroup}
\textbf{Status: Proved.} Root and fhi independently audited the complete argument.
This section treats the remaining case in D4o. All bases in the
deformation calculation are local Artinian $W(k)$-algebras with residue
field an algebraically closed field $k$ of characteristic two; all
deformations and isomorphisms are marked on the closed fibre. Let
$E_0$ be supersingular, let $F_0$ be ordinary, and choose
$0\ne P\in F_0[2](k)$. Put
\[
 A_0=E_0\times F_0,\qquad \sigma_0=(-1,1),\qquad
 \tau_0=t_{(0,P)}\sigma_0,\qquad X_0=A_0/\langle\tau_0\rangle.
\]
The conclusions are
\[
 R_{X_0}^{\mathrm{hull}}
   =W(k)[[t,s,u,v]]/(s^4-ts^2-2s),                       \tag{S1}
\]
and flatness of $\operatorname{Pic}^{\tau}$ for every Artinian
deformation of $X_0$. The ring in (S1) is finite free of rank four over
$W(k)[[t,u,v]]$. Together with D4o, this gives the two cases of the full
Igusa calculation; the ordinary ring is
$W(k)[[t,s,u,v]]/(s^2+2s)$.

\paragraph{Step 1: the cover and the multiplicative subgroup.}
As in D4o Step1, every deformation $X/R$ uniquely lifts the marked
finite etale cover $A_0\to X_0$ and its deck involution. A chosen lift
of its origin makes $A/R$ an abelian scheme; write its deck involution
as $\tau=t_c\sigma$, where $\sigma$ is a group automorphism lifting
$\sigma_0$. Thus $\sigma^2=1$ and $(1+\sigma)c=0$.
The sources are Stacks tag04DY and Mumford--Fogarty,
Proposition6.15, pp.124--126, already opened in D4o.

Write $G=A[2^\infty]$ and $D=G^0$. The connected--etale sequence,
and Cartier duality applied to that sequence for $G^\vee$, give
\[
 G/D\simeq (\mathbb Q_2/\mathbb Z_2)_R,
 \qquad M:=((G^\vee)^{\mathrm{et}})^\vee
       \simeq\mu_{2^\infty}\ \subset D.                 \tag{S2}
\]
The marked special fibre identifies $M$ with the multiplicative
part of $F_0[2^\infty]$. These are canonical subgroups: the only
etale part of $G_0$ and of $G_0^\vee$ comes from $F_0$.
The quotient $G/M$ is a Barsotti--Tate group, being Cartier dual to
$(G^\vee)^0$. Its connected part is $D/M$, so $D/M$ is a connected
Barsotti--Tate group deforming $E_0[2^\infty]$.
For these facts one can use Tate, \S1.3--1.4, pp.160--161,
\S2.2, Proposition1 and equation(4), pp.162--164, and
\S2.3, p.166 \cite{DGJ_Tate1967}; the exact sequence and duality
are also obtained level by level from their finite flat versions.

The action of $\sigma$ on $M$ is $+1$, and its action on $D/M$ is
$-1$. Here no eigenspace splitting in characteristic two is being
assumed. A homomorphism between Barsotti--Tate groups over $R$ that
reduces to zero is zero: its image is killed by a fixed power of two
in the infinitesimal formal kernel, whereas its source is two-divisible.
This is the rigidity argument in Katz, Lemmas1.1.1--1.1.3,
pp.141--143, using the equivalence between connected Barsotti--Tate
groups and divisible formal groups. Applying it to $\sigma-1$ on
$M$ and $\sigma+1$ on $D/M$ proves the assertions.

\paragraph{Step 2: the canonical divisible system of actual points.}
Choose a compatible generator of $T_2F_0(k)$, or equivalently points
$e_r\in F_0[2^r](k)$ for $r\geq1$, with $e_1=P$ and
$2e_{r+1}=e_r$; put $e_0=0$. The formal group
$\widehat A(\mathfrak m_R)$ is killed by $2^N$ for sufficiently large
$N$ (Katz, Lemma1.1.1). For $r\geq0$ define
\[
 j(2^{-r})=[2^N]\widetilde e_{N+r}\in A(R),             \tag{S3}
\]
where $\widetilde e_{N+r}$ is any lift to $A(R)$, and take $N$ large.
Such lifts exist by smoothness. Changing the lift or increasing $N$
does not change (S3): the difference of two choices, or of
$2\widetilde e_{N+r+1}$ and $\widetilde e_{N+r}$, belongs to the
killed formal kernel. These identities also give
$2j(2^{-r-1})=j(2^{-r})$. Thus (S3) defines a functorial additive map
\[
 j:\mathbb Z[1/2]\longrightarrow G(R).
\]
Its values really are two-power torsion: their reductions are torsion
and a further uniform power of two kills the formal kernel. Every
value of $j$ is fixed by $\sigma$, since the chosen special-fibre
points are fixed and the same independence argument applies to their
lifts. Set $c_0=j(1/2)$ and $\kappa=j(1)=2c_0\in D(R)$.

Write $c=c_0+z$, where $z\in\widehat A(\mathfrak m_R)=D(R)$.
The last equality follows either from the formal-group description
of $D$ or directly by killing the formal kernel by a power of two.
The affine involution condition is
\[
 0=(1+\sigma)c=\kappa+(1+\sigma)z.
\]
Since $\sigma=-1$ on $D/M$, this forces
\[
                         \kappa\in M(R).                \tag{S4}
\]
This is the step that uses the existence of the affine deck map, not
merely a lift of its linear part.

\paragraph{Step 3: construction and uniqueness of the elliptic subgroup.}
As fppf sheaves over $R$ there is a presentation
\[
 G\simeq
 (D\oplus\underline{\mathbb Z[1/2]})/
       \{(-n\kappa,n):n\in\underline{\mathbb Z}\},       \tag{S5}
\]
whose map to $G$ is $(d,r)\mapsto d+j(r)$. This is a statement on
arbitrary test schemes $T/R$, not just on $R$-points. The image of
an element of $G(T)$ in the constant etale quotient is locally a
constant rational class in $\mathbb Z[1/2]/\mathbb Z$; choose a
representative $r$ locally and subtract the base section $j(r)_T$.
The difference lies in $D(T)$. The kernel consists exactly of the
pairs $(-n\kappa,n)$, which proves (S5) as a sheaf quotient.

By (S4) the analogous subquotient
\[
 H=(M\oplus\underline{\mathbb Z[1/2]})/
       \{(-n\kappa,n):n\in\underline{\mathbb Z}\}
       \ \subset G                                      \tag{S6}
\]
is defined. It fits in an exact sequence
$0\to M\to H\to\mathbb Q_2/\mathbb Z_2\to0$.
It is a Barsotti--Tate group, as can be checked without invoking an
extension-closure assertion: multiplication by two is surjective on
both end terms and therefore on $H$; for every $n$ the snake lemma
and $M/2^nM=0$ give
\[
 0\longrightarrow M[2^n]\longrightarrow H[2^n]
 \longrightarrow (\mathbb Z/2^n)_R\longrightarrow0.     \tag{S7}
\]
Thus $H[2^n]$ is a torsor under the finite flat group $M[2^n]$ over
the constant finite etale scheme $(\mathbb Z/2^n)_R$; it is finite
locally free of rank $2^{2n}$. These finite levels give a height-two,
dimension-one ordinary Barsotti--Tate group. Their monomorphisms into
$G[2^n]$ are closed immersions. Its special fibre is exactly the
given subgroup $F_0[2^\infty]\subset A_0[2^\infty]$.

Serre--Tate, in the general form of Katz Theorem1.2.1,
pp.143--144, now produces an elliptic scheme $F/R$ and a
homomorphism $F\to A$ lifting $F_0\hookrightarrow A_0$ and
realizing $H\hookrightarrow G$. This map is a closed immersion:
it is proper and quasi-finite, since its only geometric fibre is a
closed immersion, hence finite; the map on structure sheaves is
surjective by reduction and Nakayama. We identify $F$ with its image.
Both $M$ and $j$ are fixed by $\sigma$, so $\sigma|_F=1$ by
Serre--Tate full faithfulness. The quotient $E=A/F$ is an elliptic
scheme, and $\sigma$ induces $-1$ on it by rigidity.

This subgroup is unique among lifts of $F_0$ fixed pointwise by
$\sigma$. Indeed such a lift $F'$ has multiplicative subgroup $M$
and the same etale quotient. In (S3) one may choose all the lifts
inside the smooth scheme $F'$, so every $j(r)$ belongs to $F'$.
Consequently its Barsotti--Tate subgroup contains $H$; (S7) and
the ranks at every finite level force equality. Full faithfulness
then gives $F'=F$. In particular $F$, $E$, and the subsequent
classification are functorial in the deformation.

\paragraph{Step 4: an elementary graph description of $(A,\sigma)$.}
Let $\rho:A\to E=A/F$. The homomorphism $1-\sigma$ kills $F$, so
it factors uniquely as $a\rho$, where $a:E\to A$. Since
$\sigma=-1$ on $E$, one has $\rho a=[2]_E$.
The homomorphism
\[
 \Phi:E\times F\longrightarrow A,\qquad(e,f)\longmapsto a(e)+f
\]
is an fppf epimorphism with kernel
\[
 K=\{(e,h(e)):e\in E[2]\},\qquad
 h=-a|_{E[2]}:E[2]\longrightarrow F[2].                 \tag{S8}
\]
Indeed one can lift $\rho(x)$ by dividing it by two fppf locally in
$E$, and then the residual difference lies in $F$; the kernel formula
is immediate. Thus $\Phi$ is an isogeny of degree four.
On the special fibre $a(e)=(2e,0)$, so $h$ reduces to zero.
Its composite with $F[2]^{\mathrm{et}}\simeq\mathbb Z/2$ is
zero, by invariance of maps into an etale scheme under a nilpotent
immersion. Hence $h$ factors through $F[2]^0\simeq\mu_2$.
Cartier duality gives
\[
 \operatorname{Hom}(E[2],\mu_2)
  =\operatorname{Hom}(\mathbb Z/2,E^\vee[2])
  =E^\vee[2](R).                                       \tag{S9}
\]
Let $Q$ denote the corresponding point; it reduces to zero.

Conversely, elliptic deformations $E,F$ and such a point $Q$ define
$h_Q$ by (S9) and the abelian quotient
\[
 A_Q=(E\times F)/\operatorname{graph}(h_Q).               \tag{S10}
\]
The involution $(-1,+1)$ preserves the graph, because $-1=+1$ on
$E[2]$, and descends to $\sigma$. The factor $F$ embeds in $A_Q$,
and the quotient by it is $E$ via $[2]$ on the first factor.
The marked special-fibre identification is induced by
$(e,f)\mapsto(2e,f)$. The homomorphism $a$ is $e\mapsto[(e,0)]$,
so this construction recovers exactly (S8) and $Q$. By Step3 the
elliptic subgroup used in the reconstruction is unique. We have
therefore classified all the required $(A,\sigma)$, not just those
which already come with an elliptic fibration.

\paragraph{Step 5: the affine norm equation removes the $F$ parameter.}
The marked deformation space of $E_0$ has a parameter $t$.
Choose a coordinate $s$ at zero on the dual elliptic deformation,
and let $g_E(t,s)=0$ be its two-torsion equation. The ordinary
elliptic deformation $F$ has Serre--Tate parameter $q\in1+\mathfrak m$.
Thus the pairs $(A,\sigma)$ of Step4 have complete local base
\[
 B=W(k)[[t,q-1,s]]/(g_E(t,s)).                           \tag{S11}
\]
The marking identifies the formal group of $F$ with
$\widehat{\mathbb G}_m$. The canonical point $c_0$ in Step2 lies
in $F$, by its construction there, and Katz's pushout recovery,
pp.150--152, gives $[2]c_0=q$ in this formal-group coordinate.

Choose a formal-variety section $b(u)$ of
$\widehat A\to\widehat E$ over $B$, through the identity; this
exists by formal smoothness and a choice of relative formal
coordinates. It is not required to be a homomorphism. Every formal
point of $A$ can then be written uniquely as $b(u)+(1+v)$, with
$1+v$ in the formal group of $F$. Since $\sigma=-1$ on $E$,
write $(1+\sigma)b(u)=n(t,q,s,u)$ in the formal group of $F$.
It is a principal unit and satisfies $n(t,q,s,0)=1$ identically.
The affine condition becomes exactly
\[
                 q\,n(t,q,s,u)(1+v)^2=1.               \tag{S12}
\]
The derivative of its left side minus one with respect to $q$ is
one at the closed origin. The formal implicit-function theorem
therefore eliminates $q$ uniquely. This elimination is valid over
the possibly nonreduced base (S11): lift $n-1$, which is divisible
by $u$, to a power series before imposing $g_E$, eliminate $q$
there, and then impose $g_E(t,s)$, which is independent of $q$.
The functor of triples $(A,\sigma,c)$ is consequently represented
by
\[
                    W(k)[[t,s,u,v]]/(g_E(t,s)).          \tag{S13}
\]

Forgetting the chosen origin is formally smooth onto
$\operatorname{Def}_{X_0}$, by the cover-lifting and smooth-origin
argument in D4o. It is tangent bijective: a marked isomorphism of
the covers is a translation by a formal point $d$, followed by a
marked group isomorphism. The group isomorphism preserves the
classified data; the translation changes $c$ by $(1-\sigma)d$.
On $k[\epsilon]/(\epsilon^2)$ this last expression is zero,
because the derivative of $\sigma_0=(-1,1)$ is the identity in
characteristic two. Hence no tangent parameters in (S13) are
identified. This proves the hull property. As before, the assertion
is miniversality for unpointed surfaces, not representability of that
functor; changes of origin can act at higher orders.

\paragraph{Step 6: the exact supersingular equation.}
Choose $E_0:y^2+y=x^3$, which represents the unique supersingular
elliptic isomorphism class over $k$. A universal marked deformation
is
\[
                  E_t:\ y^2+txy+y=x^3,
                  \qquad\Delta=t^3-27.                 \tag{S14}
\]
For completeness this normal form follows directly from general
Weierstrass equations. At the displayed special equation, changes
$x=\lambda^2x'+r$, $y=\lambda^3y'+\lambda^2 b x'+d$ have an
invertible Jacobian on the four coefficients
$(a_2,a_3-1,a_4,a_6)$: modulo two its diagonal entries, in the
order $(r,\lambda-1,b,d)$, are $(3,-3,-1,-1)$.
Thus these four coefficients can be normalized formally to
$(0,0,0,0)$, leaving $a_1=t$. There are no infinitesimal pointed
automorphisms of an elliptic curve. This proves the asserted
one-parameter marked deformation chart. The canonical principal
polarization identifies $E_t^\vee$ with $E_t$.

Near $O=[0:1:0]$, use $s=X/Y=x/y$ and $w=Z/Y=1/y$.
The curve equation and inversion are
\[
 w+t s w+w^2=s^3,\qquad
 (s,w)\longmapsto
 \left(-\frac{s}{1+ts+w},-\frac{w}{1+ts+w}\right).       \tag{S15}
\]
There is a unique formal solution $w=w_t(s)\in(s^3)$.
Because the denominator in (S15) is a unit and $w\in(s)$, the
fixed-point ideal of inversion is generated by
$s(2+ts+w)$. The following exact identity in the curve ring removes
the auxiliary coordinate:
\[
 (w-1)(w+2+ts)=s^3-ts-2.
\]
Since $w-1$ is a unit at the origin, the two-torsion ideal is
\[
 (s(2+ts+w))=(s^4-ts^2-2s).                             \tag{S16}
\]
Equations (S13) and (S16) prove (S1). All of the supersingular
two-torsion lies in this formal neighborhood. The polynomial is
monic in $s$ and has its lower coefficients in the maximal ideal of
$W(k)[[t,u,v]]$; Weierstrass division makes the complete quotient
free with basis $1,s,s^2,s^3$. At $2=t=0$ its two-torsion factor is
$k[s]/(s^4)$, displaying the fourth-order thickening explicitly.

\paragraph{Step 7: algebraization and flatness of the torsion Picard scheme.}
The hull has an algebraic realization over its complete local base.
The elliptic deformation $E_t$ is given by (S14); the ordinary
elliptic family $F_q$ is effective using its Weierstrass deformation
chart and the Serre--Tate coordinate isomorphism. The point $Q$ is
an actual section of its finite two-torsion scheme. Hence (S10) is
an abelian scheme over the complete ring after substituting the
solution of (S12). The compatible point $c$ is effective as well:
choose an affine neighborhood of its closed-fibre point in $A$;
the compatible Artinian sections give its coordinate values in the
complete base ring. The resulting affine involution is free, because
its proper fixed locus has empty closed fibre.

There is a useful explicit ample bundle, so no unpolarized formal
existence assumption is needed. Since the kernel of
$\Phi:E\times F\to A$ lies in $(E\times F)[2]$, multiplication
by two factors as $g\Phi$ for an isogeny $g:A\to E\times F$.
Pull back the product ample bundle by $g$, then tensor it with its
$\tau$-pullback. Swapping the factors supplies a $\tau$-linearization.
It descends to an ample bundle on the finite etale quotient $X$.
This gives the claimed algebraized projective hull family.

Finally the canonical subgroup constructed in Step3 gives
\[
 0\longrightarrow E^\vee\longrightarrow A^\vee
 \longrightarrow F^\vee\longrightarrow0,
\]
with $\sigma^\vee=-1$ on the kernel and $+1$ on the quotient.
Put $G_\sigma=\ker(1-\sigma^\vee)$. Exactly as in D4o Steps4--5,
fppf division by two in $E^\vee$ and cyclic line-bundle descent give
\[
 0\longrightarrow E^\vee[2]\longrightarrow G_\sigma
 \longrightarrow F^\vee\longrightarrow0,
 \qquad
 0\longrightarrow\mu_2\longrightarrow
 \operatorname{Pic}^{\tau}_{X/R}\longrightarrow G_\sigma
 \longrightarrow0.                                    \tag{S17}
\]
Both surjections are torsors in the fppf topology. Their degrees
are four and two. In the second sequence an invariant relative
class is represented fppf locally by a line bundle; trivialize the
base factor in its invariance isomorphism, then adjoin a square root
of the cocycle unit to linearize it. The descended class is torsion
modulo the neutral component by the norm identity, just as in (I7).
These steps do not require any section of a nonsmooth torsor over
the original base. Thus $\operatorname{Pic}^{\tau}_{X/R}$ is finite
locally free of degree eight over $F^\vee$ and is flat over $R$.
This conclusion holds for every Artinian deformation already by
Steps1--3 and (S17), independently of the explicit hull calculation.

\paragraph{Finite certificate and scope.}
The code \texttt{check\_igusa\_supersingular.py} records the exact
Weierstrass discriminant, the inversion identities, the polynomial
identity in (S16), the normal-form Jacobian, and the special-fibre
length four. Its output does not purport to mechanize the
Barsotti--Tate construction, fppf descent, or the hull argument.
Those are supplied in full above. Root and fhi independently checked the subgroup construction, the
classification, the norm elimination, the hull and effectivity arguments,
and the two torsor sequences. Their audits are recorded in the
check-event ledger.
